\documentclass[border=10pt]{standalone}
\usepackage{fontspec}
\setmainfont{Noto Serif}[Ligatures=TeXOff]
\XeTeXgenerateactualtext=1
\newfontfamily\bengalifont{Noto Serif Bengali}[Script=Bengali, NFSSFamily=bnfam]
\usepackage[Latin,Bengali]{ucharclasses}
% group-free font switching (as in header-common.tex): robust inside braces, cells and quotes
\makeatletter
\def\BN@saved{\rmdefault}
\def\BN@bn{bnfam}
\newcommand\BNin{\edef\BN@tmp{\f@family}\ifx\BN@tmp\BN@bn\else\let\BN@saved\BN@tmp\fi\fontfamily{bnfam}\selectfont}
\newcommand\BNout{\fontfamily{\BN@saved}\selectfont}
\makeatother
\setTransitionsFor{Bengali}{\BNin}{\BNout}
% Typographic quotes and dashes are not in a ucharclasses class (class 0), so a Bangla word right
% after “ or ‘ would not trigger the Bangla font. Treat them as Latin so that every switch fires.
\XeTeXcharclass"2018=\XeTeXcharclass`A
\XeTeXcharclass"2019=\XeTeXcharclass`A
\XeTeXcharclass"201C=\XeTeXcharclass`A
\XeTeXcharclass"201D=\XeTeXcharclass`A
\XeTeXcharclass"2013=\XeTeXcharclass`A
\XeTeXcharclass"2014=\XeTeXcharclass`A
\XeTeXcharclass"2026=\XeTeXcharclass`A
\usepackage{newunicodechar}
\newunicodechar{।}{{\bengalifont\char"0964}}
\newunicodechar{॥}{{\bengalifont\char"0965}}

\usepackage[table]{xcolor}
\usepackage{booktabs,tabularx,array}
\definecolor{bapilrule}{HTML}{1F2937}
\definecolor{bapilstripe}{HTML}{F3F4F6}
\renewcommand{\arraystretch}{1.28}
\setlength{\tabcolsep}{5pt}

\begin{document}
\begin{minipage}{17cm}
{\large\bfseries সারণি ৪।}\ {\large বাংলাদেশের নির্বাচিত ভিত্তি-সূচক}\par\vspace{7pt}
\small
\rowcolors{2}{bapilstripe}{white}
\begin{tabularx}{\linewidth}{@{}>{\hsize=1.045\hsize\raggedright\arraybackslash}X>{\hsize=1.545\hsize\raggedright\arraybackslash}X>{\hsize=1.000\hsize\raggedright\arraybackslash}X>{\hsize=0.409\hsize\raggedright\arraybackslash}X@{}}
\arrayrulecolor{bapilrule}\toprule
\rowcolor{white}\textbf{সূচক} & \textbf{প্রতিবেদিত মান} & \textbf{BAPIL-এর জন্য তাৎপর্য} & \textbf{উৎস} \\
\midrule
জাতিসংঘ ই-সরকার উন্নয়ন সূচক (EGDI), ২০২৪ & ১৯৩ দেশের মধ্যে ১০০তম, স্কোর ০.৬৫৭০ (২০২২ সালে ১১১তম); স্বল্পোন্নত দেশগুলোর মধ্যে সর্বোচ্চ EGDI & ভিত্তি হিসেবে ডিজিটাল সেবার কাঠামো বিদ্যমান & [5, 21] \\
আইটিইউ গ্লোবাল সাইবার সিকিউরিটি ইনডেক্স, ২০২৪ & শীর্ষ স্তরে অবস্থান প্রতিবেদিত & কাঠামো আছে; চাপের মধ্যে স্থিতিস্থাপকতা এখনো যাচাই প্রয়োজন & [22, 23] \\
অক্সফোর্ড ইনসাইটস সরকারি এআই প্রস্তুতি সূচক, ২০২৩ & ৮২তম & শাসন, ডেটা ও সক্ষমতা স্তম্ভে উন্নতির সুযোগ & [24, 25] \\
বিশ্ববিদ্যালয়ে এআই কম্পিউটিং, ২০২৫ সালের অনুমান & জরিপকৃত বিভাগপ্রতি ২ থেকে ৫৫টি GPU; সরকারি বিশ্ববিদ্যালয়গুলো নিম্নপ্রান্তে & যৌথ জাতীয় কম্পিউটিং ব্যবহার; পাইলট মডেল ছোট রাখা & [25] \\
উপাত্ত সুরক্ষা ও উপাত্ত ব্যবস্থাপনা & এপ্রিল ২০২৬-এ ব্যক্তিগত উপাত্ত সুরক্ষা আইন, ২০২৬ ও জাতীয় উপাত্ত ব্যবস্থাপনা আইন, ২০২৬ পাস & প্রথমে উন্মুক্ত ডেটা নকশার স্পষ্ট আইনগত ভিত্তি & [7, 8] \\
জাতীয় এআই নীতি & খসড়া ২০২৬-২০৩০ জানুয়ারি ২০২৬-এ জন-পরামর্শে; ২০২৬-এর চতুর্থ ত্রৈমাসিক থেকে নতুন প্রণয়ন & কারিগরি উপকরণ দেওয়ার উপযুক্ত সময় & [10, 14] \\
নেটওয়ার্ক স্থিতিস্থাপকতা & জুলাই ২০২৪-এর দেশব্যাপী বিপর্যয়ে একক ব্যর্থতা-বিন্দু উন্মোচিত & অফলাইন পরিচালনা, ব্যাকআপ ও পুনরুদ্ধারের পরিকল্পনা & [23] \\
\bottomrule
\end{tabularx}
\par\vspace{5pt}{\footnotesize\itshape সূচকগুলো ভিন্ন ভিন্ন বছরে ভিন্ন বিষয় পরিমাপ করে, তাই সরাসরি তুলনীয় নয়। মানগুলো উদ্ধৃত উৎসে যেমন প্রতিবেদিত তেমনই দেওয়া হয়েছে।}
\par\vspace{3pt}{\scriptsize\color{gray} উৎস: BAPIL টেকনিক্যাল ব্লুপ্রিন্ট, সংস্করণ ১.০ (২০২৬)। CC BY 4.0।}
\end{minipage}
\end{document}
